\clearpage
\NoSourceHead\selectlanguage{english}
\pdfbookmark[0]{About this edition}{edition-notice}
\begin{center}\Large About this edition\end{center}
\noindent Carl Friedrich Gauss, \textit{Werke}, Volume X, Part 1: \textit{Nachträge zur reinen Mathematik}. Published by the Königliche Gesellschaft der Wissenschaften zu Göttingen, in commission with B. G. Teubner, Leipzig, 1917.

This partial original-language edition contains \textit{Kleinere Veröffentlichungen}: its title and four publications on printed pages 3–8, including Zimmermann's appended note and Schlesinger's historical editorial remarks. It also contains the \textit{Arithmetik} section title and Gauss's \textit{Einige asymptotische Gesetze der Zahlentheorie}, groups I–IV, with the complete historical commentary on printed pages 16–18. The letter to v. Zimmermann dated 12 Merz 1797 and its historical editorial note follow on printed pages 19–21. The quadratic-residue notes then follow through groups I–IV, including the 1805 letter to Olbers, the determination of the decident, the fifth proof of quadratic reciprocity, and their historical commentary on printed pages 22–36. Groups V–VI contain the cubic and biquadratic residue notes and the notes, tables and lattice diagram on complex integers, with their historical commentary on printed pages 37–50. Groups VII–X follow on printed pages 51–64: further biquadratic-residue notes, the general fundamental theorem, the notion of an intensor and the theory of the residue $1+i$, with their diagrams, numerical tables and Bachmann’s historical commentary. Groups XI–XII contain further residue formulas and the correspondence on cubic and biquadratic residues: the letter to Sophie Germain of 30 April 1807, two excerpts to Olbers and one to Bessel, with Schlesinger’s historical remarks, on printed pages 65–77. The notes on the theory of forms, groups I–V, include polygonal numbers, representations by three squares, ternary quadratic forms, class numbers and complex forms, with Bachmann’s and Schlesinger’s historical commentary on printed pages 78–95. The Algebra section title and five letters concerning the fundamental theorem of algebra follow: two from Johann Friedrich Pfaff to Gauss and three from Gauss to Drobisch, Schumacher and Möbius, with Schlesinger’s historical commentary on printed pages 99–110. The notes on the circle-division equation, groups I–III, follow on printed pages 111–119: periods, the golden theorem and irreducibility, with Schlesinger’s historical commentary. The correspondence on the regular seventeen-sided polygon follows on printed pages 120–126: Pfaff to Gauss in 1802, with Pfleiderer’s construction, and Gauss to Gerling in 1819, with the complete historical footnotes. The notes on coefficients and power sums follow on printed pages 127–128, then the algebraic theorem and three letters to Schumacher on printed pages 129–134, with Schlesinger’s and Alfred Loewy’s historical commentary. The Analysis title, Klein’s attribution note and the complete \textit{Exercitationes mathematicae}, with Schlesinger’s historical remarks, follow on printed pages 137–144. The early lemniscatic-function notes, subsection I on series expansions and addition theorems, follow with their complete historical remarks on printed pages 145–151. Subsection II follows on printed pages 152–159 with ten entries on double periodicity, product representations and series expansions of numerators and denominators, including the historical numerical calculations and editorial footnotes. Subsection III follows on printed pages 160–166 with seven entries on division of the lemniscate, the numerical calculations and editable diagram, and Schlesinger’s complete historical remarks on subsections II–III. Subsection IV follows on printed pages 167–171 with six entries of miscellaneous formulas and Schlesinger’s complete historical remarks. The notes on the arithmetic-geometric mean follow through subsections I–VIII on printed pages 172–216, including the historical series, numerical calculations and relations to the perimeter of an ellipse, differential and functional equations, the nine entries on functions arising from the arithmetic-geometric mean, an editable geometric diagram, and the fourteen entries of the general theory of the lemniscatic sine with both historical manuscript-variant footnotes, followed by the notes on differentiation and the five entries on algorithms for arithmetic-geometric means and the three entries on the bilinear algorithm with its historical manuscript-variant footnote. Subsection IX, entries [1]–[14], follows on printed pages 217–231 with the relations between the infinitely many arithmetic-geometric means, the numerical examples, series transformations, elliptic integrals and rectification of the ellipse. The correspondence on printed pages 232–251 contains all eight letters: Pfaff to Gauss in 1799 and 1800, with the complete appendix from H. C. Schumacher’s Gaussiana; Gauss to Bessel in 1805 on coefficient tables and continued fractions; the exchanges between Schumacher and Gauss in 1808 and 1816 concerning Pedrayes’s problem, higher transcendental functions and arithmetic-geometric means; and Gauss to Bessel in 1828, with the two appended letters between Pfaff and Gauss. The historical postscripts and footnotes are included. The editorial introduction and the complete four-section outline of the theory of the arithmetic-geometric mean follow on printed pages 251–257, including the algorithms, differential relations, series and integral representations and their historical footnotes. The complete historical explanation of subsection I follows on printed pages 257–261, including the analysis explicitly credited to J. Horn, the inverse-function discussion and the historical account of the manuscript chronology. The complete historical explanation of subsection II follows on printed pages 261–266, with the series, elliptic-integral relations, notation distinctions, historical corrections and manuscript discussion. The complete historical explanation of subsection III follows on printed pages 266–267, including the limiting argument, its nonconstant counterexample and the connection with the lemniscatic integral. The complete historical explanation of subsection IV follows on printed pages 268–274, including the Euler and Jacobi discussion, series and integral formulas, manuscript chronology and the complete historical notes. The complete historical explanation of subsection V follows on printed pages 274–278, with the lemniscatic-sine and theta-function formulas, the transformations, the historical discussion and all six notes. The complete historical explanation of subsections VI–VII follows on printed pages 278–279, with the derivative formula and the discussion of the modular function. The complete explanation of subsection VIII follows on printed pages 279–280, including its theta-function relations, historical discussion and two notes. The complete historical explanation of subsection IX follows on printed pages 280–283, including the modular-function discussion, series, ellipsoid formulas and all seven notes. The complete historical explanations of the correspondence follow on printed pages 283–286, with the power-series argument, the quotation from Gauss’s letter to Olbers, the biographical discussion, all seven notes and Schlesinger’s signature. The notes belonging to the theory of transcendental functions follow through subsection I, entries [1]–[16], on printed pages 287–307, including the series and product identities, the arithmetic-geometric mean and elliptic-integral relations, and all eighteen historical notes. Subsection II follows on printed pages 308–310 with the solution of the inversion problem for the elliptic integral of the first kind. Subsection III follows through its opening and entries [2]–[8] on printed pages 311–320, including both editable diagrams, the coefficient series, the inverse series, the potential formulas and all seven historical footnotes. The historical introduction and complete explanations of subsections I–III follow on printed pages 320–325, including the theta transformations, reference table, manuscript chronology, elliptic integrals, all five historical footnotes and Schlesinger’s signature. The series notes on $F(\alpha,\beta,\gamma,x)$ begin with the complete four articles on continued fractions, including both sets of auxiliary series, the historical manuscript-variant footnote, the theorem and transformation. Their complete historical remarks follow on printed pages 330–331, with the comparison formulas, Lambert note, manuscript chronology and Schlesinger’s signature. Section II contains all five articles and Schlesinger’s historical remarks on printed pages 332–337, with the series relations, differential transformations, continued fractions, numerical values, logarithmic derivatives and complete logarithmic calculations. Section III includes its title, all ten articles and Schlesinger’s historical remarks on printed pages 338–354, giving relations through number (46), the coefficient table, parameter shifts, differential equations and transformations, products and series, continued fractions and arithmetic-geometric mean formulas, with the historical footnotes. Section IV includes its title, all six articles and Schlesinger’s historical remark on printed pages 355–359, giving two further series relations, the trigonometric expansion and its coefficient recurrence, a transformation with its sixth-power example, differential equations for a further series and arguments about incomplete sums, with the historical notes. The correspondence includes the seven letters: Gauss to Bessel of 21 October 1810, Gauss to Olbers of 17 October 1811, Gauss to Bessel of 21 November and 18 December 1811, Gauss to Laplace of 30 January 1812, Gauss to Bessel of 5 May 1812, and Gauss to Laplace of 5 November 1812, on the logarithmic integral, transcendental and factorial functions, definitions, complex integration, logarithms, continued fractions, probability, least squares and ellipsoidal attraction, with all historical notes, on printed pages 360–379. Schlesinger’s complete remarks follow on printed pages 379–381, with quotations from Legendre and Laplace and the historical account of the manuscripts and certified copies. The following notes on series begin with the title and articles 1–8 of the method for summing a divergent series, including the differential relations, coefficient expansions, numerical calculation, entries from Leiste and the Tagebuch, and historical manuscript notes, on printed pages 382–385. Schlesinger’s complete commentary on printed pages 385–389 discusses Euler’s divergent series, asymptotic expansions and Bessel functions, with its historical quotations, references and notes. Section II, on the basic concepts of series, comprises six articles and Schlesinger’s remarks on printed pages 390–395. Section III, Questions on the metaphysics of mathematics, follows with its five entries, lattice diagram and Schlesinger’s remark on printed pages 396–397. Section IV gives six series for discontinuous functions and Schlesinger’s remark on printed pages 398–399. Section V contains its introduction, both printed articles on convergence and imaginary quantities, and the appendix on printed pages 400–406. Section VI follows with all nine articles on complex quantities, their boundaries, integration and periodic expansions, including the boundary diagram, through printed page 419. Section VII includes all five articles on the convergence of the midpoint-equation expansion, with the numerical tables and historical editorial notes, through printed page 428. The following correspondence includes the letter to C. F. Hindenburg of 8 October 1799, the letters to Schumacher of 4 and 6 December 1849, 5 February and 1 September 1850, and the letter to H. G. Grassmann of 14 December 1844, with their historical notes, on printed pages 429–437. The historical remarks on sections V and VI and the complete mathematical explanation of section VII follow on printed pages 437–442. The complete historical commentary on the correspondence follows on printed pages 442–447, including the manuscript chronology, J. Horn’s quoted calculation, all historical notes and Schlesinger’s signature. The Geometry title and the four notes of Ansätze zur transzendenten Trigonometrie follow on printed pages 451–452, with all five diagrams. Stäckel’s complete historical commentary follows on printed pages 453–456, with its four diagrams, functional equations and spherical-trigonometry formulas. The note on spherical trigonometry and Stäckel’s historical remark follow on printed pages 457–458, retaining the attribution to Wittstein’s textbook. The Geometry correspondence includes letters [1]–[7], dated 19 March–21 April 1836, on printed pages 459–467, with nine diagrams and the historical footnotes. Stäckel’s historical closing remark follows on printed page 468. The notes on manifolds of n dimensions begin with section [I], Über die Ermittelung des kleinsten Werthes für das Maass des Zwanges, on printed pages 469–472, with the quotation from A. Ritter’s 1853 dissertation and its historical footnote. Section [II] follows with parts [1] and [2] on printed pages 473–481, from the 1850/51 lecture Über die Methode der kleinsten Quadrate as recorded by August Ritter, followed by Stäckel’s historical remark on printed pages 481–482. The diary section follows with its historical title and the complete first manuscript sheet, containing the seven entries dated 30 March–24 May 1796 and both continued fractions, transcribed as editable text and drawings. The title of the printed diary, the complete historical prefaces by Felix Klein and Ludwig Schlesinger, and entries 1–146 with their complete historical commentary follow on source leaves 512–615, through entry 146, the closing note, the complete diary subject index, the historical correction list Klein’s concluding remarks on the volume and the complete printed contents and printer imprint. Five further manuscript entries on source leaf 495, from Methodus facilis through resolvi oporteat, ten prose entries on source leaf 496, from Aequationis through detexi, and six entries on source leaf 497, corresponding to printed entries 51–53, 55 and 57–58, together with four entries on source leaf 498 corresponding to printed entries 59–62, and four entries on source leaf 499 corresponding to printed entries 64–67, and the final entry on source leaf 500 corresponding to printed entry 71, together with eight entries on source leaf 501 corresponding to printed entries 72–75 and 77–80, and five entries on source leaf 502 corresponding to printed entries 81, 83–85 and 88, are included. Source leaf 503 is represented by the two opening prose entries, the first series corresponding to printed 91a, and the two final separate formulas corresponding to part of printed 91b; its intervening calculations remain outside this selection. Eight manuscript entries on source leaf 504, corresponding to printed entries 92–94 and 96–100, are included; entry 95 remains outside this selection. Source leaf 505 includes the six entries corresponding to printed entries 101–106; the calendar note 107 remains outside this selection. Source leaves 506–507 contain the complete manuscript entries corresponding to printed entries 108–120. Source leaf 508 includes the entries corresponding to printed entries 121 and 123–132; entry 122 remains excluded. Source leaf 509 includes the manuscript entries corresponding to printed entries 133, 135–137 and 139–140; entries 134 and 138 remain excluded. Source leaf 510 includes the manuscript entries corresponding to printed entries 141–143 and 145–146; entry 144 remains excluded. Source leaf 511 is blank and is represented by an anchor. The rest of the reproduced manuscript text on source leaves 493–510 is outside the current transcription. The source PDF scope is \GaussXOneEditionScope. Blank leaves 9, 17, 103, 105 and 143 are accounted for in the editable source without added reader blanks. The half-title on source leaf 4 and its blank verso on leaf 5 are included. The other preceding title matter and subsequent text are outside this edition.

The source is the Göttingen Digitisation Centre scan, identifier PPN236018647. The PDF's modern introductory leaf precedes the historical scan; its PDF leaf numbers are one greater than the corresponding IIIF canvas numbers.

\noindent\href{https://gdz.sub.uni-goettingen.de/id/PPN236018647}{Read the historical source at the Göttingen Digitisation Centre.}

A second scan of the 1917 volume, held by the Hong Kong University of Science and Technology Library, supplies the legible margins of the Germain, Pfaff and analysis correspondence and is compared for printed pages 65–482 and 485–577, the unnumbered correction leaf and the complete printed contents and printer imprint.

\noindent\href{https://lbezone.hkust.edu.hk/rse/?p=13258}{Historical Gauss volumes at HKUST Library.}

The text is AI-transcribed and AI-typeset. Gauss's writing and the historical contributions of Johann Friedrich Pfaff, E. A. W. Zimmermann, Ludwig Schlesinger, Bachmann, Alfred Loewy, Felix Klein, H. C. Schumacher, J. Horn, Stäckel and Galle retain their original authorship. The cited Italian work is by Joseph Piazzi; the cited historical German translation is by Johann Heinrich Westphal. No translation of this historical text is supplied here. No human or independent review is claimed.

Historical spellings, wording, formulas and punctuation are retained without mathematical correction. In Schlesinger's quotation on printed page 4, the source visibly spaces the words as ``gegebe nhat''; that spacing is preserved without asserting the printer's intention. The second denominator series in I [4] has two terminal dots; its neighbouring series have three. In II [1], the small numerator coefficient is read as $1/8$ and the power as $4$ from the enlarged printed contours; the denominator retains $(l2)(l3)(l4)(l5)$. Modern fonts, including the bold representation of two historical blackletter inscriptions, spacing and reflow do not reproduce the identity or dimensions of historical type. All text, formulas, tables, rules and diagrams are editable LaTeX; no source-page images form this edition's body.

In the quadratic-residue notes, the repeated final factor in I [1] and the differing exponent parentheses on printed page 31 are retained. The notation in group II uses two fraction bars. On printed page 30, three small exponent contours are represented as $u$ where they differ from neighbouring $\mu$ forms; their original printing intent, including the possibility of defective type, remains uncertain. The historical notes and corrections are retained in their original positions.

In group VI, blank index cells, repeated entries and the printed signs are preserved. Detached horizontal ink on printed page 44 and small tails on certain final digits are not confidently separate mathematical signs. The lattice diagram is an editable drawing of the observed point pattern, lines and labels; its regular spacing and stroke widths do not reproduce the irregularities of the printed impression.

In group VIII, both small displayed exponents of $i$ are represented by $\rho$, although the surrounding historical prose mentions $\rho$ and $\rho'$. A distinct raised prime is not visible in the second contour; the original printing intent remains uncertain. The intensor diagrams and the $A/B$ diagram preserve the observed labels and connections in editable drawings. Group X retains the printed numerical pairs without recomputation, the nested brackets, the literal $[b]$, the underbrace and its term in their observed positions, and the crossed four-circle reference mark. Its two wide tables are reflowed in successive groups of columns. Detached ink is not treated as additional mathematical content.

In group XI, the Göttingen scan does not show a closing parenthesis after the exponent expression $\tfrac12(a+A-2$ on printed page 69; the HKUST scan shows it. The body retains the Göttingen reading without asserting the original printing intent. Group XII retains the historical French spellings, the example list ending in 94, the literal numerical equalities, the raised zero in $7^0\,6'$, the omission points and Gauss’s quoted corrections. No arithmetic has been substituted for the printed numbers.

The notes on forms preserve the open square and triangle symbols, the numerical decompositions, matrix entries, blank cells and unbarred roots. In group II, both displayed substitutions for $f$ and $g$ have positive signs in the source. The historical discussion of the depression method is Bachmann’s commentary. The forms in group V retain the differing primed factors and repeated products as printed.

The algebra correspondence retains its underlining, letterspacing, omission points and unbarred roots. The crossed-stroke sign in Pfaff’s reference to Newton is represented by \# without expanding the historical abbreviation. The letters and their historical editorial notes retain their distinct authorship.

The circle-division notes retain the two empty root-expression cells, unbarred roots and nested exponents. Products printed without parentheses remain unparenthesized. The integral-shaped abbreviation is retained in the text together with the historical explanation.

The seventeen-sided-polygon construction is an editable drawing of the printed lines, arcs and labels. Its stroke dimensions represent the observed construction without claiming historical printing identity. The letters and the editor’s quoted account of Huguenin retain their distinct authorship.

The power-sum formulas retain their printed coefficients and signs, including the negative term involving $C$ in the general expression. The algebraic note retains the unbarred root and raised zeros. The letters retain their omission points and the historical statement that a letter’s ending is missing. Loewy’s discussion is historical editorial commentary; its quoted wording, inequalities and open question remain part of the original text.

The \textit{Exercitationes} retain their recursive coefficient notation, unbarred roots, unparenthesized products and literal omission points. The three unlabelled figures are editable drawings of the observed strokes and connections. In the product on printed page 144, the Göttingen scan lacks the opening parenthesis visible in the HKUST copy; the principal reading is retained. A detached mark between the plus-minus sign and the number on printed page 142 has a low contour in the principal and an extended vertical contour in the counterpart; it is not confidently a separate mathematical sign. Small coefficient digits touching fraction bars remain qualified by their enlarged contours. Klein’s attribution correction and Schlesinger’s remarks are historical editorial writing, not AI-authored commentary.

The early lemniscatic notes preserve the unbarred roots, blank numerical cells and intermediate values. In entry [3], the small exponent has an upper 2 and an ink-filled lower contour read as 3; the representation $2/3$ remains qualified by those contours. In entry [4], the Göttingen scan lacks a closing parenthesis visible in the HKUST copy, whose thin column rules are also absent from the principal. Entry [7] lacks a visible opening bracket before its page locator in the principal. Its third right-hand root contains two adjacent stacked digit columns, 2 over 2 and 4 over 5, with an interruption between their horizontal strokes; the positions are retained without assigning a multiplication meaning or resolving the original compositorial intent. Schlesinger’s historical corrections remain in his remarks and footnotes.

In subsection II, the cubic contours in entry [1] vary from open to ink-filled and 8-like; 3 is a qualified reading, including the exponent on 1,766. The printed 0,8 is retained without recomputation. The two initial product rows in entry [2] have no intervening fraction bar. The historical corrections to the coefficients of $x^{12}$ remain in their footnotes, without replacement of the earlier formulas. The numerical stacks retain their omitted signs and blank cells; entry [9] retains the closing square bracket after the decimal prefix $0,$. Entry [8] has no visible opening bracket before S. 67; its repeated $\lambda\lambda$ exponent is retained. The roots are unbarred where printed without an overbar; the root in the historical footnote defining $\varpi$ has an overbar. Detached ink is not assigned an additional linguistic meaning.

Subsection III preserves the coefficient arrays, blank positions, three hollow squares, repeated index factors and literal numerical calculations. The lemniscate diagram is an editable drawing with two crossing lobes and four transverse marks. Its geometry and stroke dimensions represent the observed figure without claiming identity with the printed strokes. The small denominator in the exponent $1/4$ in entry [7] is ink-filled in the principal and clearer in the counterpart; the small fourth-power contours retain this qualification. Historical editorial corrections remain in Schlesinger’s remarks and footnotes.

Subsection IV retains the incomplete expressions, absent signs, unlabelled numerical values and differing uses of $\pi$ and $\varpi$. The powers omitted from two coefficient numerators remain absent; Schlesinger’s discussion is retained as historical commentary. Small fourth-power contours are ink-filled in the principal and clearer in the counterpart. The HKUST copy has a gathering mark 22-star on printed page 171 that is not visible in the principal.

The arithmetic-geometric-mean notes preserve the exponent shaped as $u$, the stars and blank terms, the differing root bars, leading and trailing primes, the overbar over $1+\nu$, and the unmatched closing brace in II [40]. The three small cubic powers in II [7]–[9] include ink-filled and 8-like contours; 3 remains a qualified reading. The small fourth powers and fractional exponents retain their observed positions. The Göttingen page references S. 48 and S. 56 have only low fragments before S.; the complete opening brackets visible in the HKUST copy are not supplied in the principal body. S. 83 lacks a visible opening bracket in both copies. Detached marks are not assigned additional mathematical or linguistic meaning. The printed numbers and the two historical manuscript-variant footnotes remain without recomputation or correction.

Subsections III–IV retain the unclosed parentheses, the absence of a terminal equality after the first Hinc expression, the ten literal points before (1), the unbarred fourth root, the repeated label [24], and the absent [23]. The printed denominators 36384 and 8024 and the entry 0,8472130.8 are retained. In IV [6], the principal has a gap between the final denominator factors 16 and 20 where HKUST shows a point; no point is supplied. The second exponent in the initial definition of $t$ on printed page 192 has a 6-like contour in the principal and a closed 8 in HKUST; 8 is retained as a qualified reading. Small cubic contours elsewhere include ink-filled forms. The final preceding identity has a square on its first coefficient, while the following proof begins with a cube. The triangular figure retains its lines, hatching and five labels as editable geometry; modern stroke dimensions do not claim historical type identity.


Subsection V retains the fourteen entries, the unclosed source citation and the unmatched brackets, the separate leading primes, the differing barred and unbarred roots, and the literal numerical omission points. On printed page 197, the period notation has two separate serif vertical strokes without a connecting top bar in both scans; it is represented by $\mathrm{II}$ without identifying its intended mathematical role. The low point after $\mu$ in the opening differential has no assigned multiplication meaning. In [16], the first inner closing parenthesis survives as a low fragment; its following outer parenthesis and terminal period are retained. HKUST has a comma at that terminal position. Formula [17] retains its final $\sin\varphi$ and its differing bracket position. The historical two-row calculation and both manuscript-variant footnotes retain their printed signs.

The small cubic and fifth-power contours in V include ink-filled forms. In the final $W\varphi$ product, the last two displayed denominators are represented as $a^3$; the final principal contour supports a qualified 3 reading, whereas the HKUST contour resembles 5. The differing numerator powers and the two positive exponential powers in [13] remain unchanged. These qualified glyph readings do not establish the original printing intent. Detached ink and damaged parenthesis fragments are not assigned additional primes or operators. The two scans and modern reflow do not establish historical type identity.


Subsection VI retains $M(xy)$ without a supplied comma and the unchanged derivative numerators where the denominators carry primes. Its final differential expression has two stacked rows without a joining sign or fraction bar; the unfinished parentheses remain unfinished. The damaged closing parenthesis in the first differentiation formula is complete in the HKUST copy. Detached low ink beneath the first $dx$ is not confidently linguistic. Subsection VII retains the nested unbarred roots, the terminal minus after $2ff'$, the term $zz^{9/2}$ and the distinct differential notations. The small fourth power on $b$ in the $2ff'$ expression is ink-filled in both scans; 4 is a qualified reading. Several small cubic contours, including the denominator $f^3$, are also ink-filled or 8-like; 3 remains a qualified reading. The original printing intent is not established by these representations.


Subsection VIII retains the two leading primes in the opening recurrence and the denominator $4x$ in [4]. The first sine-root series on printed page 215 has two visible omission points in its denominator in the principal scan; the HKUST copy has a faint third middle point. Its small numerator exponent 9 is ink-filled and 3-like in the principal, and clearer in HKUST. Small cubic contours also include ink-filled forms. On printed page 216, the first numerator root lacks an opening parenthesis while retaining its closing parenthesis in both scans. The main expression and its historical manuscript-variant footnote retain their differing signs. The footnote's raised terminal point has no multiplication meaning. Detached marks beside the $x^4$ expression and [5] are not confidently additional linguistic or mathematical content. Damaged title and citation brackets retain the forms supported by the counterpart; qualified glyph readings do not establish original printing intent.


Subsection IX, entries [1]–[9], retains the literal numerical tables and their blank cells, the appended half-fractions, the descending six-fraction staircase and the different root bars. The short horizontal stroke above the leading zero of $0{,}2100421$ on printed page 219 is represented by an overbar; it is absent from the HKUST copy, and no negative-value interpretation or original printing intention is asserted. The small cubic contours in the proof and its historical manuscript-variant footnote include ink-filled and 8-like forms; 3 remains a qualified reading. The bracketed exponent $[-1]$, the final root $\surd(\delta+\gamma M)$ and the repeated $A''A''$ remain as printed. Raised terminal points carry no multiplication meaning. Detached ink is not assigned additional mathematical content. The historical editor's footnote and manuscript citations retain their original authorship.


The concluding entries of subsection IX retain the unbarred roots, the literal bracketed tangent operators, the repeated factors and the differing coefficient groupings in the ellipse formulas. The numerical examples and omission points are transcribed as printed. Raised punctuation has no multiplication meaning. The closing rule after printed page 231 in the HKUST copy is absent from the principal and is not supplied in its body. Detached ink above the first quotient on that page is not assigned additional mathematical content.


In the analysis correspondence, the omission points, differing coefficient forms and Pfaff’s historical correction footnote are retained. The second letter’s left-margin losses in the Göttingen scan are supplied by the HKUST copy. A detached mark after the comma following “verliehren” is absent from that copy and is not confidently linguistic. The four root indices at the start of printed page 237 retain their short underlying strokes; the radicands have no overbars. Raised terminal punctuation has no multiplication meaning.


In the letter to Bessel of 3 September 1805, the fourth defining fraction retains the small denominator power read as $f^3$; its closed contour leaves the original printing intent uncertain. Weak closing-parenthesis contours in the first continued fraction are supported by the clearer HKUST copy. The bold coefficients and letters correspond to the red manuscript characters identified by the historical footnote. The printed May date in the 1808 letter and its historical April correction remain separate, as do the body’s $xx+uu=1$ and the footnote’s manuscript reading $xx+yy=1$.


The four-section historical outline retains the repeated equation number (8), the square on $c$ in (4), the second $k^2$ term in (19), and $x^{10}$ in (35). The unequal signs have vertical strokes. The weak closing parentheses in (29) and after the references to (33) and (12a) are supported by the HKUST copy. The eighth point of the last omission row in (5)$\prime$ is raised in both scans; its position supplies no additional operator. Detached ink is not assigned mathematical or linguistic content.


The explanation of subsection I retains the gamma-function argument $8/6$, the repeated $xx$, the starred epsilon exponents and the positive fourth-power term in the long footnote. Its two unstarred constants are epsilon, distinct from the Latin $e$ in the gamma-function formulas. The Nielsen reference reads ``l. (1)'' in both copies; no preceding letter is supplied. Small cubic contours and damaged parentheses retain qualified glyph identity. The period after 365, visible in the HKUST copy but absent from the principal, is not supplied. Detached ink has no assigned linguistic or mathematical meaning.


The explanation of subsection II retains $663$, $41/16048$, the reference to III, page 261, and the absent factor $a$ in the second term of the displayed formula for $q$. The two-prime reference on printed page 265 and the missing closing parenthesis after the reference on printed page 264 are preserved. Bold and italic $L$ and $M$ follow the historical editor’s explicit distinctions. The faint numerator $9$ in the coefficient for beta and weak parenthesis contours retain qualified glyph identity; the second scan supplies the closing parenthesis in the denominator of $[29]'$. Detached ink is not assigned extra mathematical meaning.


In the explanation of subsection III, the principal reference to page 149, article [8.], has a damaged closing square bracket and no visible outer closing parenthesis before the comma; the latter, visible in the HKUST copy, is not supplied. Weak parentheses elsewhere retain their qualified readings from the two copies. Both unbarred roots, the overbar on $c$ and the exponent $2\pi i/\log2$ are preserved.


In the explanation of subsection IV, the principal Euler series has two visible points in each of two omission runs, where the HKUST copy has three. The principal lacks a visible opening square bracket at the K-prime note marker and an opening parenthesis before the page-188 reference; the closing forms remain. In the first series of the faculty-prize note, the principal has a gap between the numerator factors 3 and 5, where HKUST has a point. These fuller counterpart forms are not supplied. The damaged beta, small cubic exponents and broken punctuation elsewhere retain qualified readings. The date 14. Dezember 1798 and the historical editor’s numerical corrections remain as printed. Fraktur and italic z remain distinct.

In the explanation of subsection V, the principal first theta-series argument has no visible parentheses; the second lacks its opening parenthesis and argument separator. The fourth series has a gap where the HKUST copy shows a minus, and the last W relation lacks a closing argument parenthesis. The principal also omits the separator in the first doubled-argument expression in the following prose. These gaps are retained. In the cited title on printed page 277, the principal reads \textit{Summa io}; the fuller t contour in HKUST is not supplied. Broken parentheses, small powers and the low closing-parenthesis fragment in the squared Q expression retain qualified identities. The differing signs in the exponential denominators, barred and unbarred roots, and historical corrections remain unchanged.

The explanation of subsections VI–VII retains the repeated products, all prime strokes, the two terminal points in its heading and the decimal $0{,}1660958116$. The weak closing-parenthesis contours in the reciprocal-mean formulas remain qualified. Detached ink near the derivative is not assigned a new symbol. The historical editor’s correction remains in its original discussion.

The explanation of subsection VIII retains the identical lower indices 01 in the numerator and denominator of its relation $[5]^{\prime\prime}$. The absent opening parenthesis before \textit{vergl.}, the missing closing parenthesis in the first theta argument of $[5]^{\prime}$, the missing opening parenthesis in its second theta argument, and the absent argument separator in the first right-hand theta expression on printed page 280 are retained. Other weak parenthesis and exponent contours have qualified readings; fuller forms in the comparison scan do not replace the principal readings. Incidental detached ink is not assigned additional symbols.

The explanation of subsection IX distinguishes Fraktur and italic $M$, and upright and italic $a,b$. It retains $b=1-a$, both occurrences of $a^{\prime\prime}$ in the final four-relation note, the repeated factors and the printed positions of the squares after $T^{\prime\prime}$. The inline root on printed page 282 and the root in the integral of the following note have no overbars; the former has no opening parenthesis. The damaged numerator letter in the third relation of that note is read as $b^{\prime}$ with qualified confidence: both scans show an open lower contour and a taller left stroke, and the surrounding notation supports the reading without establishing the original printing intent. Weak parentheses elsewhere remain qualified. The two raised terminal points carry no multiplication meaning.

The correspondence commentary retains the historical correction of Pfaff’s coefficients, all repeated factors and indices, the spelling \textit{angentium} in the Latin title, and the unaccented Greek quotation. The low points in the denominators and the raised terminal punctuation are retained in their respective positions; the latter has no multiplication meaning. The letter quotation and the continuing historical note retain their complete wording.

In the transcendental-function notes, entry [1] retains the missing opening parenthesis in the left exponent of [1] on printed page 288, the differing manuscript variants in the historical notes, and the second note marker without a closing parenthesis. On printed page 289, the small denominator glyph following $1+$ is represented as $\xi^{\prime}$: both scans show short horizontal contours and a descending lower stroke, distinct from the neighbouring $\mu$ forms. Its glyph identity and original printing intent remain qualified; contextual expectation does not replace the printed contour. The unbarred roots, repeated factors and positions of the squares are retained. The raised terminal point has no multiplication meaning.

In entry [2], the third numerator of $P$ begins with a qualified reading $x^{2n}$, and the first expansion of $R$ has the denominator exponent $n-2$. The repeated products and visible gaps between factors are preserved. In entry [3], a gap in the principal denominator on printed page 292 is retained without supplying the minus visible in the comparison copy. Small powers and damaged parentheses have qualified readings.

The heavier function brackets in entry [4] follow the distinction stated in the historical footnote. Its weak parenthesis contours retain qualified readings. The crescent-shaped reference marks in entries [4]–[5] are editable outlines representing the observed form without assigning an astronomical meaning or reproducing historical type. Entry [5] retains the final $y^9$, distinct from the preceding $x^9$ terms, and the historical note's manuscript exponent $9/2$ remains separate from the printed formula's $9/4$. Literal omission points, repeated factors and differing signs are preserved without recomputation.

Entries [6]–[9] preserve the observed repeated factors and differing series signs. The principal copy has incomplete parentheses in the product fraction of entry [6] and lacks an opening parenthesis in the denominator of entry [7]; the fuller forms in the comparison copy are not supplied. Other visible parenthesis fragments retain qualified readings. Entry [6] has two separated omission points at the end of its opening display, where the comparison copy has an additional faint point. The small eighth-power exponents in entry [9] retain the reading $1/8$ from both enlarged witnesses. Its historical manuscript-variant note has three visible points in the principal copy and four in the comparison copy. The unbarred root and the raised terminal punctuation are retained; the latter has no multiplication meaning. These forms are transcribed without recomputing the identities.

Entry [10] retains two absent opening parentheses in the principal product on printed page 297. Its small exponent in the second factor has the qualified reading $12$ in the principal, against $13$ in the comparison copy. The historical note on printed page 298 retains its two rows of factors without introducing fraction bars or parentheses. Its $1-x^{10}$ and the identity’s quarter-power exponent $25/4$ follow the enlarged sources. The historical notes and their different markers remain separate from the body formulas.

Entry [11] preserves the two-row arrangements inside the braces, the repeated-letter products, the different squared and cubed bracket factors, and the literal five-point leaders. Its small exponent in the $t$ expansion reads $52$ in both enlarged copies. The comparison copy’s 38-star footer is absent from the principal and is not inserted into its body. Worn parenthesis contours remain qualified; modern reflow does not reproduce the original type.

Entry [12] retains $x^{87}$ in the printed formula and the historical note reporting $x^{67}$ in the manuscript. In the second expansion after “Nun ist ferner”, the principal has a gap followed by two visible points after $x^{24}$; the comparison copy also shows a faint preceding point. The principal’s visible sequence is retained. The two-row arrangements, repeated-letter factors and different powers of bracketed expressions remain as printed.

Entry [13] preserves the nine unbarred root signs, their indices, the repeated-letter factors and the printed numerical values. Its five parallel calculation blocks are arranged in two groups, with the original numerical sequences, blank positions and rules retained. The small parameter typeform is represented by $\alpha$; a closely similar italic $a$ remains a qualified alternative. The distinct flat-topped glyph in $\pi/\varpi$ is kept distinct from $\omega$. Raised terminal points carry no multiplication meaning. Both historical notes, including the note printed below entry [14]’s opening, belong to entry [13].

Entry [14] retains the final $y^6$ denominator in its opening product, repeated-letter factors, seven unbarred roots and the three distinct historical notes. The damaged exponent in $(1+\varepsilon^3)n$ on printed page 305 is read as $3$, supported by the clearer HKUST copy. The opening parenthesis before $\varepsilon+\varepsilon^4$ in that page’s final left member is not established in the principal image and is left absent; the HKUST copy shows it. In the final historical note, the square-root bar covers both $q/p$ and $P/Q$. The separate omission line contains twelve literal points. Raised terminal points carry no multiplication meaning.

Entry [15] preserves $48b^4$ in the coefficient of $[B]$ and the repeated-letter products. In the definition of $B$, the principal image shows one omission point in the numerator and two in the denominator; the HKUST copy shows three in each position. Only the visible principal points are represented. The fragmented exponent $6$ in the second polynomial is supported by the clearer HKUST copy. These differences do not imply a mathematical emendation.

Entry [16] retains both printed $T^3$ terms, including $6ttT^3$, without changing the displayed polynomial to agree with the following factored expression. Its two fourth roots have no overbars; the printed positions of the tangent and sine powers, the repeated products and the equation label are retained.

Subsection II retains the printed $9/2$ exponent in the opening series for $R$, the differing coefficient signs, the historical question mark and four unbarred roots. On printed page 308 the Göttingen scan shows two omission points where the HKUST scan shows three, and no visible period before the citation’s closing bracket where HKUST has one; the body follows the Göttingen impression. Two damaged exponent numerators are represented as $25$ with support from the second scan, without asserting the printing intent of the weakened contours.

In subsection III the manuscript-deletion markers and the historical note remain with the text. The editable circle drawings retain intersection and containment, including the smaller circle’s displacement within the larger one. The coefficient series retain their unfinished parentheses and omission stars. The Göttingen impression shows a low round point after $16786k^{10}$ which is absent from HKUST; it is retained without multiplication meaning or certainty about printing intent. The damaged parenthesis in the definition of $k$ and several small power contours are read with the second scan. The first manuscript-page reference retains the Göttingen form [S 172].

Entry III [4] retains both plus signs at the historical line join in the formula for $\log t$, the unfinished parentheses and terminal signs in the coefficient series, and the star above the omission points. The final $\log t$ series has two visible omission points in the Göttingen impression and three in HKUST; the body retains two. The final $u^9$ in the preceding series has a damaged principal contour corroborated by the second scan. The dated manuscript annotation and its historical footnote remain unchanged.

Entry III [5] retains the positive initial $u/U$ in the logarithm series and the later negative cosine terms. The printed $B^{\prime\prime\prime}\cos3E$ in [11] remains distinct from the sine term in [10]. The trailing plus sign, $2/2$, repeated-letter products and historical note are retained without mathematical correction. Small cubic powers and a fragmented parenthesis retain their qualified readings from the two scans; detached ink is not supplied as additional content.

Entry III [6] preserves both $1/8$ factors in the summation exponents, the repeated-letter products and the three unbarred roots. The first denominator in [18] has no closing parenthesis before its second root; the damaged remaining parentheses retain their observed positions. The small parameter is represented by $\alpha$; a closely resembling italic $a$ remains a qualified alternative. The fragmented opening bracket of S. 224 and the small powers retain their source qualifications. The historical \textit{bedarf Berichtigung} annotation and its footnote remain alongside the printed formula, without retrospective alteration.

Entry III [7] retains the differing signs and repeated factors in its expressions for $B$, including the positive initial term of $B/A^3$. The small cubic contours and the final digit of the exponent $13$ are ink-filled in the two scans; the reading retains its qualification. The obscured final letter of \textit{einfach} in the Göttingen copy is legible in HKUST. The historical annotation and manuscript-variant note retain their authorship and wording.

Entry III [8] retains its revised formulas and both historical notes, including the editor’s statement about the integral sign absent from the manuscript. The signs, repeated factors, unbarred roots and differing products remain as printed; the historical corrections are not applied retrospectively to preceding entries. Small ink-filled cubic contours retain the earlier qualification.

The historical commentary retains the page range 436--415 in its note about Hattendorff’s copy. Its account of manuscript comparison belongs to the historical editor. In the first theta-function expression for $P(x,y)$ on printed page 321, the argument divider has a short raised vertical contour; an editable stroke preserves that position without extending it or identifying it as a prime. The later expressions have full-height dividers.

The commentary on printed page 324 retains the ratio $P/Q$ beside the theta ratio with indices $00/10$, and its differing asterisk counts in the text and footnote. On printed page 325, the inner parenthesis in $(r(\lambda))^2$ is a qualified reading of a fragmentary principal contour, supported by the other scan. The first theta argument divider has only a short low vertical contour. The following numerator in the definition of $k$ has no established divider in the principal; the other scan shows a broken bar. The final integral retains the absent opening parenthesis before $1-\tau^2$. These readings are not algebraic emendations.

The continued-fraction notes preserve repeated products and the printed coefficient $5\mathord{.}10\mathord{.}12$ in $T$. The historical footnote gives $\frac59s^4$ within the last term, whereas the body has $\frac39s^6$. On printed page 328, two principal ellipses show only two separated points, and one denominator has $13\ 15$ without a visible product point. The other scan shows a third point in the final left-hand series. The detached horizontal segment at the start of the first continued-fraction rule is treated as part of the interrupted rule; its printing intent remains qualified.

In article [3.], the third nonconstant coefficient of $Q$ retains the $2$-shaped numerator digit visible in both scans. Its exact printing intent remains qualified; the surrounding pattern has not been substituted. Product points absent in the principal are left absent, including gaps within several coefficient strings and the theorem in article [4.]. The other scan supplies clearer points in several of those positions. The transformation ends with two visible principal points in each continued fraction, whereas the other scan has three. Both roots in article [3.] have no overbar.

\par The historical remarks on printed pages 330–331 retain the word form ``autet'' without a supplied initial letter. In the first $G$ expression, the Göttingen scan has no punctuation between $2$ and $\frac{13}{2}$; the comparison copy shows a low point at that position. The closing parenthesis of the inline $F$ expression on printed page 331 is a qualified reading of its surviving upper and lower fragments, supported by the second scan. The three terminal points after the four division relations include a weak final contour. The comparison copy’s 42-star gathering mark is absent from the Göttingen impression and is not inserted into its body.

\par In section II, articles [1]–[3], damaged braces and parentheses are retained as qualified readings, supported where legible by the comparison copy. Two small cubic exponents on printed page 333 are ink-filled; their reading as $3$ remains qualified. The transformed equation retains the positive second term and the unsquared factor $16(1-t)$. The last denominator of the general continued fraction remains $C.C+\gamma$. The specialized fraction has four diagonal omission points and two separated points in the principal right-hand denominator; the comparison copy has three in the latter position. A detached segment aligned with a fraction rule is not assigned an additional operator meaning. The complete historical transformation note is retained.

\par Article [4] preserves the printed numerical values and their historical qualification. Several small denominators in the $F$ arguments touch the fraction bars and have $3$-like contours in the Göttingen scan; the reading as $5$ is qualified, with clearer corresponding stems in the comparison copy. The three cubic denominator exponents in $\Lambda^{\prime\prime\prime}z$ are heavily inked and remain qualified. In the fourth-order term on printed page 335, the first denominator has no established opening parenthesis in the principal and is retained as $z+1)^4$; the comparison copy has the opening parenthesis. Thirteen literal omission points follow the expansion. Detached spots and a protrusion adjoining a parenthesis are not assigned additional mathematical meanings.

\par Article [5] retains the printed calculation stacks, blank digit positions and differing numerical endings, including $93424$, $94762$ and $94586$. No values have been recalculated. The bracket in the opening $2\log_{10}e$ line closes immediately after $0,$. The roots have no overbars, including the fourth root of $8$. The integral denominator’s opening parenthesis is a qualified reading of its broken principal contour, supported by the comparison copy. Interrupted calculation rules are represented as rules; their detached segments receive no additional operator meaning. The nine-row numerical table retains the row labelled only $0{,}75$, the final $564$ in the first row, and eighteen omission points.

\par The historical remarks to section II preserve the first reference $[3]$ without a terminal point, the later $[3.]$, the reference to article $[4.]$ for the generalization, and »Zur der Formel«. The principal has no visible point after »vergl« before »oben S.~145«; the comparison copy has one. The fourth root carries an overbar; the square root of $\pi$ does not. The incomplete gathering siglum is represented by its two visible diagonal strokes, without assigning a separate letter identity. Their modern geometry does not reproduce the historical type.

\par Section III begins with the printed repeated products, four terminal omission points and the nine numbered relations. The first coefficient numerator on printed page 339 remains $\alpha+r.\beta+r$ without added parentheses; the later two occurrences of $x-xx$ remain unchanged. The point after $(\alpha+r)/\alpha$ is raised and has no multiplication meaning. The damaged closing parenthesis in the initial derivative fraction and the opening parenthesis before $\alpha+r$ in the combined coefficient remain qualified readings of the contours in both copies. The historical wording, including »Hingegen werden sich«, »Werthe« and »seyn«, is retained.

\par Article [2] retains the six-row coefficient table and its three blank coefficient cells. In the fourth row, the right-hand numerator begins $Cr+\alpha)$: neither copy supplies an opening parenthesis between $C$ and $r$. The denominator retains the qualified, partly broken opening parenthesis in $\alpha(r+\gamma)$. Equation (26) has a prime on both occurrences of $F$ in both copies; both are retained. The repeated $\gamma\gamma$ in (13), $x-xx$ and the order $\beta+\alpha$ in (23) and (24) remain unchanged. Fragmentary parentheses elsewhere are read from their surviving curved strokes and the corresponding positions in the other copy; their exact printing condition remains uncertain. Detached spots beside the equations and the counterpart table are not assigned additional mathematical meaning. The complete starred historical note and its literal opening bracket followed by a comma are retained.

\par Article [3] retains the four differential equations for $P$ and $Q$, the leading minus in the differential quotient, the exponent $-\gamma$, the repeated factors in the product for $A$ and the printed positive value $1/\pi$. The two omission groups contain three points each. The terminal point after $1/\pi$ is raised near the fraction bar and has no multiplication meaning. The final series keeps its unparenthesized repeated factors and both literal denominator products. The partly broken inner parenthesis in the first equation for $P$ and the final small $\beta$ retain their qualified contour readings; the other copy provides corresponding curves and the lower part of $\beta$. The detached mark below »allgemein« in the principal and the counterpart spot below the product are not additional punctuation.

\par Article [4] retains the three distinct continued-fraction endings: equation (39) ends with a denominator $1-\mathrm{etc.}$; equation (40) places »etc.« below the last coefficient without another fraction bar; equation (41) places it after the last $x$. The products retain their unparenthesized factors and low points, and equation (41) keeps $xx$ beside $x^3$. The partly broken argument parentheses in the principal are read with the corresponding curves in the other copy; their upper fragments are not additional primes. The historical bracketed footnote and its separate equation remain part of the text.

\par Article [5] retains the unbarred root, the four literal coefficient series, $ff$ beside $f^2$, $A^{(0)}$ beside $A^0$, and the Roman index IV. No equation label (42) is visible beside the four series in either copy, and none is supplied. The point above the comma after $fA^{(\gamma-1)}$ in the principal is represented by a semicolon-shaped pair; the other copy shows a comma alone. This preserves the principal contour without settling its printing intent. The continued fraction (46) ends with the printed $1+\mathrm{etc.}$. Both bracketed historical notes and the X 1./44 signature remain in the text.

\par Article [6] retains eleven unbarred roots, the expressions $\operatorname{tang}\theta^2$ and $\cos\theta^2$, the repeated products $\xi\xi$, $\zeta\zeta$, $\xi^\prime\xi^\prime$ and $ff$, and the Roman indices I, II, IX and X. The two numerical continued fractions retain their distinct coefficients and endings. The reference (42) remains as printed. The closing historical note gives the angle with a raised zero and one prime stroke; its manuscript calculation is not supplied by the printed passage.

\par Article [7] distinguishes the Fraktur constants $\mathfrak{A},\mathfrak{B}$ and their primes from italic $A,B$, and retains the repeated products $xx$, $tt$ and $\mu\mu$. Its three raised terminal points are punctuation without multiplication meaning. The partly broken digit in $2\mu$ and weakened parentheses retain qualified readings from both scans; the short parenthesis fragments beside $F$ are not represented as extra primes. The bracketed label [31a], the letterspacing of \textls[120]{erstlich} and \textls[120]{Zweitens}, and the historical note concerning later additions remain in the text.

\par Article [8] retains the six Roman labels, the repeated $xx$, the six unbarred roots and the two occurrences of $\mathfrak{A}$ in relation V. Relation IV has a point at the end of its first line and another at the beginning of the continuation in both scans; both positions are preserved. Its exponent $\gamma-\beta-1$ and final parameter $i+\gamma$ remain unchanged. The definition of $\Pi z$ retains the integral without displayed limits and gives its limits in the prose; the product has three literal omission points after 4. Weak parenthesis contours remain qualified readings of the two scans.

\par Article [9] retains the three coefficient rows, the repeated $xx$, the unreduced fractions $2/2$ and $2/3$, and the three primes in the expression for $R$. Its factor sequences and low product points are preserved without added grouping parentheses. The points after the first definitions of $A$ and $B$ are raised punctuation without multiplication meaning. Small fraction digits and thin parenthesis contours remain qualified readings of the scans.

\par Article [10] uses Latin $a$ and retains its five relations for $\Pi a$, including the three-point omission sequences, the product points, the three primes and the bracketed second condition $[\text{pro }f=\infty]$. The derivative expansion has no terminal punctuation after that bracket. The final operator is printed as cotang; the closing rule precedes the historical remarks.

\par Schlesinger’s historical remarks retain their three paragraphs, both bibliographical notes, the signature and the printed references. After \emph{dürften}, the principal scan has a comma-shaped mark; the counterpart has an additional upper point. After the reference S. 221, a period is visible in the counterpart but absent from the principal image. The text follows the principal at both positions without asserting the original printing intent. The root in $b_1=\sqrt{ab}$ has a bar. Detached ink below an equality, near a parameter and in the margins is not treated as an additional mathematical sign. The historical editorial discussion is preserved without using it to alter Gauss’s text.

\par The opening of section IV retains the bracketed title, the May 1809 source description and the first two articles. In the principal image, points are absent within the spaced abbreviation $\text{u s w}$ and after 1809 in the source description, and between the first two factors of the second series numerator; corresponding points are visible in the other scan. The principal’s visible spacing is retained. Both copies have a low point after $\alpha+1$ in the first relation’s final $F$ expression; it is preserved without an assigned multiplication meaning. The coefficient display contains nineteen literal omission points, and its two shorter omissions contain three points each. Repeated $ff$, all indices, the bracketed word [zweite], the historical note and the raised-star printer’s signature remain unchanged. Broken parenthesis contours and the last digit of the reference S. 289 remain qualified readings of the two scans.

\par Article 3 retains the bracketed denominator $n[+1]$, three omission points and a terminal point after the opening minus sign, the complete sixth-power expansion and its unparenthesized factors. The final bracketed factor $[{.}\,f^6]$ and the absence of a product point between the two denominator $\Pi$ factors are preserved. Thin or interrupted parenthesis contours remain qualified by comparison of the two scans. The historical footnote’s raised terminal point has no assigned multiplication meaning. Modern line breaks do not reproduce the historical layout.

\par Article 4 retains the printed unbalanced parenthesis sequence $((\alpha+1)(\beta+1)-(\gamma+1)\delta+1))$, visible in both scans; no opening parenthesis is supplied before $\delta+1$. The repeated products $xx$ and $tt$, the $ddP$ numerators and the squared $x$ and $t$ in the differential denominators remain unchanged. The final raised terminal point is punctuation without an assigned multiplication meaning. A detached mark beside the other scan’s first second-derivative fraction is not treated as an additional operator.

\par Article 5 preserves the two three-point omissions in the quotient for $S$ and the two five-point omissions between $d$, $g$ and $p$, without supplying signs. Its two paragraphs, the four displayed sign comparisons, the three-line prose list and the bracketed manuscript note remain distinct. Historical spellings, including allemahl, nehmlich, Summirung and nachweiset, are retained.

\par Article 6 retains $1-\theta$ in the displayed $B_\mu$ expression and $1+\theta$ in the following inline $F$ expression, as printed in both copies. The indices in the two occurrences of $b_{\mu+1}$ are represented in lowered type; their thin contours and shallow historical lowering remain qualified by the enlarged scans. The principal image has a gap between the first two numerator factors $k\ k-1$ in the concluding series, where the other copy shows a point; the principal reading is retained. Schlesinger’s complete signed remark, its historical dating statements and its citations remain part of the historical text.

\par In the letter to Bessel of 21 October 1810, no period is supplied after ``wie Sie, where both scans show a gap. The staggered series retain their trailing signs, repeated letter products, braces and all omission points. The small cubic exponents follow the principal contours; the other scan has closed, eight-like forms in some positions, especially the final inline denominator. These glyph readings remain qualified.

\par The letter to Olbers of 17 October 1811 retains ``unexplicabeln, the repeated $xx$, $ddy/dx^2$ and all literal omission points. The seventh point in the opening line is a narrow vertical contour in both scans. A small upper mark immediately before the citation’s Roman II remains uncertain and is not assigned a separate linguistic sign. The HKUST copy alone shows the footer ``X1. and ``46; these are not present on the principal leaf.

\par The letter to Bessel of 21 November 1811 retains the two raised factorial notations $a^{bIc}$ and $1^{xI1}$, with the final 1 at the same raised baseline as $xI$. The serifed separating contours are represented by italic $I$ without assigning them a multiplication meaning. Small cubic contours and closely touching fraction bars remain qualified. All historical brackets, ten notes and the distinct sine expressions are retained. The complete second scan supplies letters clipped at the principal scan’s left margin in the da Cunha and Göttingische Anzeigen notes.

\par The letter of 18 December 1811 retains the differing dagger markers, the unbarred root, the raised commas, the absent punctuation after S. 444 and after anhängt, and the printed numerical sequence with 18 above its first zero. The closing parenthesis after $1{,}0+0{,}001i$ survives as separated arcs in the principal scan and more clearly in the second scan. The historical numerical correction remains in its editorial note; it is not applied to the letter.

\par The letter to Laplace of 30 January 1812 retains its historical French spellings, bracketed insertions, continued fractions and both twenty-seven-point omission lines. The opening parenthesis before $n$ in the limiting relation survives as two separate arcs in the Göttingen scan and is clearer in the second scan; it is not an additional prime. The historical note’s differing year remains in the note without changing the letter.

\par The letter to Bessel of 5 May 1812 retains its two unbarred roots, raised zeros, numerical values and historical editorial notes. The first logarithmic series has no plus sign before etc. The Göttingen scan shows no terminal point after zeigt in the Cauchy note; the point in the second scan is not added to the principal text. The raised terminal point in the manuscript-variant note is punctuation, not multiplication.





The opening notes on series retain the bracketed equality in the title and the five differential relations, including repeated $xx$ and the Roman derivative indices IV and V. The Göttingen citation has no visible period after 46a; the second scan has one. A small low mark after the second differential relation occurs only in the second scan; no terminal point is added to the principal text. Detached ink above the first $\xi$ is not treated as an accent.

The numerical calculation on printed page 383 retains its blank digit positions and the point after 125. Article [6] has no visible period in its principal heading, although the second scan has one. Its integral quotient retains the printed colons; the raised terminal point is punctuation without multiplication meaning. The historical manuscript correction remains in its note.

Articles [7.]–[8.] retain the unbarred root signs, repeated letter products and unreduced numerical factors. The historical note to [17] preserves the point and baseline 3 following its first root sign. The entries retain their printed brackets and the date Brunsv[igae,] Oct. 16. [1797].

In Schlesinger’s commentary, the principal witness has two omission points in the Mascheroni relation and the generating series for the coefficients \(\gamma\). Its date reads “16 Oktober 1797”; the other scan has a point after 16. The missing opening parenthesis before “siehe a. a. O. S. 426)” is retained. The raised terminal point in \((\beta)\) is punctuation. The historical commentary and its mathematical assessments are Schlesinger’s.

In section II, the principal scan has a point after “unterscheidet” before “alle”, and after “vorgestellt” before “das”. The other scan shows a comma at the latter position. The printed sequence of prime numbers begins with 1. These forms and the incomplete punctuation in “u s. f.” are retained.

In article 3, the reading “wäre” uses the clearer second scan where the principal has a low fragment at r. The fourth example retains “Z B” and the absent point before the comma in “u. s. f ,” from the principal impression. The wording “eine untere Grenzen”, “folge” and “so wohl” remains as printed.







The diagram in Questions on the metaphysics of mathematics retains its sixteen open points and twenty-four connecting segments in editable form. Its drawn contours and dimensions do not reproduce the historical printed impression.

The discontinuous-function series retain the unreduced products and fractions, repeated letters and raised zeros of the source. The three elevated omission points and following lower terminal point in the second formula remain distinct from its denominator points.



The historical note after the second article on convergence states that two further manuscript articles were not printed. The present scope comprises the text of the printed edition.































The conclusion of the explanation of section VII preserves the eleven-row account of the numerical table, both paired inequality signs, the repeated factors and all printed logarithms and angle values. No values have been recalculated. The four small cosine cubes have clearer open contours in the Göttingen scan and more closed contours in the comparison copy; the reading as 3 remains qualified. The detached spot below the equation label on printed page 441 is not assigned punctuation or operator meaning.

In J. Horn’s quoted calculation on printed page 446, the opening parenthesis in $F(z)$ and the parentheses around $z^2+1$ in its denominator are interrupted or absent in the Göttingen scan and legible in the Hong Kong copy. Their transcription follows the latter witness. The distinct printed forms tang, cotang and cotg, the Fraktur $\mathfrak E$, the prime on $\Phi$, the order $\vartheta i$ and the signed integration limits are retained. The historical account and its stated comparisons with manuscripts belong to Schlesinger and the named contributors.

The five diagrams accompanying the four opening geometry notes are encoded as editable lines, curves and labels. Their incidence, seven separate interior strokes in the third group and all printed labels are retained; modern dimensions and fonts do not reproduce the original drawing or type. The repeated factors, bracketed beta and superscript comma in $\psi^{0,1}$ are retained. Detached spots beside the denominator in group II and between or below the diagrams are not established as additional linguistic or mathematical signs.

Stäckel’s historical commentary on printed pages 453–456 retains its two asterisks, literal omission points, repeated products and the placements of the powers in the Arctang and final cosine expressions. The second $BF$ formula on printed page 456 contains $\psi(\Delta,\tfrac12))$ as printed; no missing argument is supplied. The opening parentheses beside $\varphi$ on page 454 and $\psi$ on page 456 have detached upper contours in the Göttingen scan; comparison with the Hong Kong copy does not establish separate primes. The first paragraph’s interrupted opening parenthesis is legible in that copy. The alpha coefficients, including the subscripted forms in equation (14), retain the qualification attached to their close alpha/a contours. The raised point after equation (5) is terminal punctuation. Incidental strokes and specks are not assigned additional mathematical meaning. The four diagrams and the curved angle signs are editable drawings; their modern dimensions do not establish historical type identity. Stäckel’s explanations and statements about the manuscripts retain their historical authorship.

In Schumacher’s letter of 7 April 1836, several parentheses in the coordinate equations are interrupted in the Göttingen scan; the HKUST witness supports the parenthesis readings. The opening parenthesis in the first squared trigonometric denominator of \(U\) is likewise faint in Göttingen and clearer in HKUST. These are qualified readings of the printed forms. The roots have no overbars, and powers retain their positions after the trigonometric arguments.

In the letter of 13 April 1836, no punctuation is visible after “aufhellen lassen” in the Göttingen witness; the faint low point in HKUST is not supplied in the text. In the letter of 19 April, “Punct,” follows the comma visible in both witnesses; a detached upper speck in Göttingen is not treated as a semicolon.

The diary manuscript is represented by editable text, underlines, marginal strokes and continued-fraction layouts. Modern type and dimensions do not reproduce Gauss’s handwriting. The first sheet retains the raised insertion of “non,” its abbreviations and the visible endings of both series. The Göttingen witness has one clear point after the first series; the Hong Kong copy shows additional faint points. No trailing plus or ellipsis is supplied after the second series or its continued fraction. The initial letter of “septendecim,” some abbreviation points and the terminal ligatures are qualified handwritten readings. An upper fragment beside the point before 15 does not securely establish a colon; it is not assigned a second punctuation or division sign. The two three-stroke marginal marks receive no letter or operator identification. The title’s footer follows the Göttingen initials without supplied points. Incidental strokes, detached ink and differences in rule continuity are not assigned additional linguistic meaning.

The comma after “Hinsichtlich” in Klein’s preface is retained from both printed witnesses. Detached upper specks near “der Art” do not establish an additional quotation mark. Klein’s account of the earlier diary edition’s expansions, corrections and omissions, and Schlesinger’s account of the contributors, retain their historical wording and authorship. These accounts do not authorize changing the separately transcribed manuscript.

The Ceres and Pallas signs in the diary’s dated table are represented by editable curves and strokes. Modern type, stroke dimensions and formula reflow do not reproduce the historical printing. The faint closing parenthesis after c minus b in the fourth denominator of the continued fraction in the commentary to entry 7 remains a qualified reading of both witnesses. The formula printed in entry 12 and Bachmann’s later proposed correction are retained separately in their historical positions.

In the printed note containing the deleted entry 17a, the numerator exponent in \(m^2/n^2\) is congested in the Göttingen witness; the Hong Kong copy supports the reading 2. The editors’ transcription of that deleted entry remains distinct from the unresolved manuscript cancellation. The printed “applicata” and its note reporting manuscript “applicatum,” and the separate “qui” and “quae” statements, retain their historical positions. Barred and unbarred roots, the elevated “(p)” in entry 22 and the period outside “[30]” are retained.

Entry 32 retains its printed coefficient 185 and the historical editor’s correction to 165 separately. Its first cubic exponent is heavily inked in the Göttingen witness and clearer in the Hong Kong copy. The raised point after the footnote’s product fraction is terminal punctuation. The differing denominator forms in entry 33 and its commentary are retained without inserted parentheses or algebraic rewriting. Detached ink below a variable or above an integral sign is not assigned an additional sign.

The broken closing bracket in entry 38 is legible in the Hong Kong witness. Entry 39 retains its repeated final squared difference and the distinct original wording and historical editorial correction. Entry 49 retains the differing single-asterisk reference and double-asterisk note marker. The unbarred roots of entries 50–53 remain distinct from the barred root in the commentary to entry 50.

The diagram in diary entry 81 preserves the five printed labels, solid and dashed connections and three short curved marks as editable geometry. Additional fine marks in the Göttingen impression are absent from the Hong Kong copy and are not confidently identified as printed labels or relations. Entry 82 retains the root followed by a baseline point and \(3x\); the historical editor’s correction remains in its commentary.

In the commentary on diary entry 83, the unnumbered fraction for \(ix_i^m\) on printed page 528 retains the single visible closing curve after its final index in the Göttingen impression. The Hong Kong copy shows two curves there. Damaged type or imaging cannot be excluded. The broken closing parenthesis after \(x\) in equation (18) remains identifiable. The exponent \(m+2\) in equation (7), congested in the Göttingen copy, follows the clearer second digit in the Hong Kong copy.

In the commentary on diary entry 87, the small curved stroke to the left of the infinity limit on printed page 531 appears only in the Göttingen copy. It is not confidently a minus sign, and no mathematical meaning is assigned to it. Repeated factors and the distinct lower and middle product points in the series are retained.

In the commentary on diary entry 92, the letters o in Produkte and u in nur on printed page 535 have broken contours in both copies. These readings use the surviving contours and the linguistic context; the scans do not establish intact letterforms.

In diary entry 97, the first formulas III and IV retain the printed numerator \(\sin\delta+p\) without supplied parentheses. The Moon crescents and Aquarius sign are editable drawings of the observed symbols. A detached low mark after “ihnen.” in the Göttingen scan is absent from the other copy and is not confidently a separate linguistic sign. Galle’s concluding account and the historical credit to Chompré remain part of the historical text.

Entry 98’s historical explanation reads “st nämlich” in both scanned copies; no initial letter is supplied. Its final integral relation has a raised terminal point, represented as punctuation without multiplication meaning. A small detached mark beside \(\varpi\) before an equality in the Göttingen copy is not treated as an additional sign. The quoted numerical value and the historical note about a printing error are retained without recalculation or retrospective emendation.

In entry 101’s displayed integral, the mark after \(r^2\) in the first factor has only a lower curved stroke in both witnesses. Its representation as a damaged closing parenthesis is a qualified reading, not a claim that the complete printed form is established. On printed page 545 the HKUST copy has the gathering marks “X1.” and “69”, which are not visible in the principal scan and are not added to its historical body.

In der Erläuterung zu Tagebucheintrag 108 ist die öffnende Klammer vor \(1+\mu^2x^2\) in beiden Druckzeugen nur unterbrochen erhalten. Die Klammerlesung bleibt entsprechend qualifiziert; die moderne Type stellt den beschädigten Druck nicht nach.

In the final formula of the commentary to diary entry 113, the opening parenthesis before \(x\) survives only as a low fragment in the principal scan; the HKUST copy shows the complete parenthesis. The column divider on printed page 553 is also much clearer in the HKUST copy. The editable parenthesis and divider represent these readings without claiming to reproduce the damaged printing.

A low isolated point between Mense and Octobr. in diary entry 121 appears in the principal scan but not in the Hong Kong copy. It is not confidently identifiable as punctuation and is not added to the text.

The terminal mark after Stäckel in the commentary to diary entry 143 is comma-shaped in the principal scan. The heavier mark in the Hong Kong copy does not establish the original printing intent; a comma is retained.

The HKUST copy has the gathering marks \textnormal{X\textsubscript{1}.} and \textnormal{73} below Klein’s concluding remarks (printed page 577); these are not visible on the corresponding principal leaf. They are recorded here as a witness variant outside the principal body. Klein’s historical account of the edition and its contributors is retained as his text.

On printed contents page 579, the fourth leader point in the row referring to page 52 stands above the other points. Its raised position is retained without assigning it a mathematical meaning; the modern type does not reproduce the original printing geometry.

On printed contents page 581, the twelfth leader position in the row referring to page 280 has a comma-shaped contour in both witnesses. It is represented by a comma without assigning it a mathematical meaning or determining the intended printing.

On printed contents page 583, the principal has a separate low point between 1836 and April in the letter dated April 21. The HKUST witness has no corresponding point. The point is retained without assigning it a mathematical meaning; an incidental ink mark remains an alternative to intended punctuation.

On printed contents page 584, the principal has twelve visible leader points in the row for entry 27, with a wider interval after the seventh. The HKUST witness has thirteen points. The principal sequence and wider interval are retained; the modern spacing does not reproduce the original printing dimensions.

On printed contents page 585, the second leader position for entry 93 has a comma-shaped contour in both witnesses, represented without mathematical meaning. The principal’s final digit in the date 1799 December 23 is ink-filled; the clearer HKUST reading supports 3 without establishing the principal type’s printing condition.

On printed contents page 586, the tenth leader position for entry 146 has a comma-shaped contour in both witnesses. It is represented by a comma without assigning it a mathematical meaning.

The half-title preserves the smaller baseline 1 in BAND X 1. The comparison copy has the same title and a blank verso; its library watermark and show-through are not historical printed text.

The five manuscript entries on source leaf 495 preserve their abbreviations, the unbarred root, six prime strokes and two three-point omissions. The date Sept. 14 remains beside the Methodus row without reassignment. The final digit of Sept. 16 is heavily inked and retains a qualified reading. The deleted ending ibus before es in gaudentibus is represented as editable crossed text. Comparison also uses the original Mathematisches Tagebuch, Cod. Ms. Gauß Math. 48 Cim., leaf 2v, held by the Göttingen library (\href{https://gdz.sub.uni-goettingen.de/id/DE-611-HS-3382323}{digital manuscript}). Modern type, crossing lines and monochrome underlines do not reproduce the handwriting or ink. Detached flecks and the narrow strokes beside elicere are not assigned additional linguistic meaning.

The ten manuscript prose entries on source leaf 496 preserve the forms multiplicati aggregati, demonta and La Grangiani, the date abbreviations, underlines and two round pointed markers. The original diary leaf 3r supplies an additional witness. The open zero contour in the first equation and the heavily inked ending of secundum retain qualified readings. A low point beside the second denominator is represented without multiplication meaning. Detached flecks between the following prose entries are not assigned additional linguistic meaning. The final integral group remains outside this selection because its crossed-out matter is not securely legible. Modern type and reflow do not reproduce the handwriting or ink.

On manuscript leaf 497 the crossed word elasticam and the replacement lemniscatam are retained. The two strokes above o in cöpi are represented by an umlaut, with the precise linguistic intent remaining qualified; the printed coepi is not substituted. The raised interlinear p, the two raised-zero ordinals, the manuscript wording ad polygona p laterum, the four points after the series and the complete continued fraction are preserved. The handwritten fourth power and the looped date digit 6 retain qualified readings. Entries corresponding to printed 54 and 56 remain outside the selection because the crowded numerator and crossed-out word are not securely resolved. Comparison includes the original diary leaf 3v; modern type, underlines and reflow do not reproduce the handwriting or ink.

The four manuscript entries on source leaf 498 retain the a-shaped exponent, the qualified beta root index, the unbarred roots, the right parenthesis without a supplied opening partner, the raised (0-1), and the sum with numerator mm+6mn+nn. The crossed word plures is represented as editable text; its heavier final strokes are corroborated by the historical printed note. The date abbreviations and double-underlined ordinal ending remain qualified handwriting readings. Two short points and two longer marks after partes are represented without subtraction meaning. The three-by-three marginal stroke pattern receives no assigned mathematical meaning. The final entry is excluded because its cancellation before pollicetur is unresolved. The original diary leaf 4r is an additional witness; modern type and strokes do not reproduce handwriting or ink.

The four manuscript entries on source leaf 499 preserve the square mark, the qualified a-shaped glyph, the order of the signs and the point after +1 without multiplication meaning. The form utranque, the abbreviation ind., the date forms and the additional emphasis under puras are retained. The original diary leaf 4v shows Demonstrationes without a separate textual cancellation; the darkened reproduced strokes are not expanded into deleted letters. The following entries remain excluded because their cancellations and continuation have not been fully established. Modern type and monochrome underlines do not reproduce the handwriting or ink.

The final manuscript entry on source leaf 500 retains the abbreviation aeq., its equation, the date Jul. 27. Gott, both underlines and the adjacent stroke pattern. The pattern is represented by editable paths without assigning it a symbol identity. Comparison also uses original diary leaf 5r. Modern type and stroke geometry do not reproduce the handwriting or ink; detached flecks are not assigned additional linguistic meaning.

The eight manuscript entries on source leaf 501 retain their abbreviations, dates, underlining and three marginal stroke patterns. The crossed word before rationalem is read fract, with moderate confidence because the crossing obscures the upper joins; it is not expanded. The second occurrence following congruentias is retained as secund: the final letter has a returning hook, but separate letters um are not established. The abbreviation ending in d with an upper stroke remains unexpanded. The low point after Quomodo and the final colon after 1799 are retained. Comparison also uses original diary leaf 5v. Modern type and editable strokes do not reproduce the handwriting or ink.

The five manuscript entries on source leaf 502 preserve their abbreviations and dates. The iterated-function formula retains nine prime strokes, the unbarred cube root and the final plus without a supplied ellipsis. Its first 1 and denominator i are heavily inked and retain qualified readings. The manuscript has quavis, unlike the printed quovis; the two-part punctuation after ordine is represented as a semicolon. Both marginal stroke patterns and the three underlines remain editable. Comparison also uses original diary leaf 6r. Modern type and stroke geometry do not reproduce the handwriting or ink; faint detached marks are not assigned additional linguistic meaning.

Source leaf 503 is included selectively: its two opening prose entries, the first series and the two final separate formulas. The middle calculations remain outside this selection because their cancellations and overwritten figures are unresolved. The lowered terminal form in attractionem is read as m; the retouched radicand in the unbarred root is read as 6. Both retain qualified glyph confidence. Literal omission points, the missing closing parenthesis and the final incomplete fifth-power expression are retained without the bracketed supplements of the printed transcription. Comparison also uses original diary leaf 6v. Modern type and reflow do not reproduce the handwriting or ink.

Source leaf 504 contains eight selected manuscript entries, corresponding to printed entries 92–94 and 96–100. Entry 95 is excluded because the crossed date remains unresolved. The small year 1799 retains qualified digit confidence. The divided Jul, underlining, connecting stroke, marginal stroke patterns and unbarred root are represented in editable form. Detached spots are not assigned mathematical or linguistic meaning. Comparison also uses original diary leaf 7r; modern type and strokes do not reproduce the handwriting.

Source leaf 505 contains the six manuscript entries corresponding to printed entries 101–106. The deleted Qu and final m in the opening entry retain qualified readings, as do the retraced first A of the trigonometric series, the date numeral 6 and the continuation rime. Seriem, the unbarred root, repeated letter products, unexpanded abbreviations and marginal stroke patterns are retained. Entry 107 remains excluded because its cancellation and faint final annotation are unresolved. Comparison also uses original diary leaf 7v. Modern strokes and fonts do not reproduce the handwriting or ink.

Source leaf 506 contains six complete manuscript entries corresponding to printed entries 108–113, compared also with original diary leaf 9r. The crossed-out omnes, term, duo and raised ut are retained with qualified readings where the ink obscures letter contours. The retouched date 12, unexpanded dates, seven marginal patterns and five separate omission strokes are preserved. Modern fonts and editable paths represent the content and positions without reproducing the handwriting or ink.

Source leaf 507 contains the seven manuscript entries corresponding to printed entries 114–120. The unexpanded dett., raised ex, semicolon between the first two rho symbols, cancelled moti, marginal figure, four unbarred roots and unexpanded dates are retained. Readings obscured by cancellation or retracing remain qualified. Comparison also uses original diary leaf 9v. Modern fonts and editable paths represent content and positions without reproducing the handwriting or ink.

Source leaf 508 includes the manuscript entries corresponding to printed entries 121 and 123–132; entry 122 remains excluded because its cancelled wording is unresolved. The two low points after eruimus, capital Elementa, cancelled orbitam, unexpanded dates and raised br with its lower bar are retained. The manuscript reading simus is qualified against the printed counterpart simul; the joined letters of Hinc and other retraced contours also retain qualification. Comparison uses original diary leaf 10r. Modern fonts and editable paths do not reproduce the handwriting or ink.

Source leaf 509 includes the manuscript entries corresponding to printed entries 133, 135–137 and 139–140. The crowded opening-word ending in entry 134 and cancelled or altered mathematics in entry 138 remain unresolved, and these entries are excluded. The raised in and insertion stroke, underlined tres, unexpanded dates and abbreviations, X−1=0, repeated letter products and marginal stars are retained. Retouched letters, the evolutae ending and faint punctuation retain qualified readings. Original diary leaf 10v provides an additional witness; its extra terminal marks after 24 are not established in the principal reproduction. Modern fonts and paths do not reproduce the handwriting or ink.

Source leaf 510 includes the manuscript entries corresponding to printed entries 141–143 and 145–146. Entry 144 remains excluded because its cancelled word is unresolved. The raised essentialis and insertion stroke, four underline runs, repeated letter products, single equality in the congruence statement, unexpanded dates and Gott are retained. The compressed Seeberg. and retouched b contours remain qualified readings. The short mark above operae is not confidently a separate linguistic sign. Original diary leaf 12r supplies an additional witness; terminal marks visible only there are not supplied to the principal reproduction. Modern fonts and editable paths do not reproduce handwriting or ink.
